\section{方法与理论}
\label{sec:method}

\subsection{Black--Rangarajan对偶与GNC-QPEP前向}

\begin{lemma}[GM的精确BR对偶]\label{lem:brdual}
对任意$\mu\ge1$、$c>0$与$e\ge0$，
\begin{equation}
\begin{aligned}
\rho_\mu(e)&=\min_{w\in[0,1]}\big\{we+\psi_\mu(w)\big\},\\
w^*(e)&=\rho_\mu'(e)=\Big(\frac{\mu c^2}{\mu c^2+e}\Big)^2,
\end{aligned}
\end{equation}
从而对任意$z$有
$\min_{w\in[0,1]^N}\sum_i\{w_ie_i(z)+\psi_\mu(w_i)\}=\sum_i\rho_\mu(e_i(z))$。
标量目标在$(0,1]$上严格凸、驻点为$w^*$，盒约束是笛卡尔积；完整证明见附录~\ref{app:proofs}。
\end{lemma}

\begin{theorem}[GNC-QPEP前向]\label{thm:gnc}
在(A1)、(A2)、(A5)下，固定$\mu\ge1$并交替执行
\begin{equation}
\begin{aligned}
w_i^{(k)}&=\rho_\mu'\big(e_i(z^{(k)})\big),\\
z^{(k+1)}&=\argmin_{z\in\mathcal{M}}\sum_iw_i^{(k)}e_i(z).
\end{aligned}
\end{equation}
则：
\begin{enumerate}
\item[(i)] $w$步是精确的，$z$步是加权QPEP。消去$\mathbf{t}$得
  $\mathbf{t}(w,\mathbf{q})=G(w)^{-1}\sum_iw_iW_i(\mathbf{y}_i-R(\mathbf{q})\mathbf{p}_i)$，其中
  $G(w)=\sum_iw_iW_i$；约化代价是关于$\mathbf{q}$的四次多项式，KKT系统有四个（或七个）变量，
  $N$仅通过$O(N)$累加进入固定规模的矩阵。
\item[(ii)] $F_\mu(z^{(k)})$非增；若每个加权极小点唯一，则迭代收敛到不动点
  $(z_\mu^*,w_\mu^*)$，其中$w_{\mu,i}^*=\rho_\mu'(e_i(z_\mu^*))$且
  $z_\mu^*\in\argmin_z\sum_iw_{\mu,i}^*e_i(z)$。
\item[(iii)] 当$\mu\downarrow1$时，每个聚点满足$w_i^*=\rho'(e_i(z^*))$且
  $\Phi(z^*)=0$，即$z^*$是$F_{\mathrm{rob}}$的驻点。
\item[(iv)] 全局性质只属于内层加权QPEP：外层GNC是非凸块坐标方法，可能停在鞍点；相干多数外点
  还可使$F_{\mathrm{rob}}$的全局极小点本身出错（第~\ref{sec:discussion}节）。
\end{enumerate}
陈述(ii)是标准的单调块下降论证；(iii)由$\rho_\mu'$的连续性与加权问题的一阶条件得到。在
3D--3D实验中，内层求解器是QPEP的闭式标量权重$4\times4$特化；$4\times4$的Davenport/Horn
变体在附录~\ref{app:systems}中做了基准测试。
\end{theorem}

\begin{algorithm}[t]
\caption{前向GNC-QPEP（一次鲁棒求解）}
\label{alg:forward}
\footnotesize
\begin{algorithmic}[1]
\Require $\{(\mathbf{p}_i,\mathbf{y}_i,W_i)\}_{i=1}^N$，尺度$c$，调度$\gamma$
\State 初始化$\mathbf{q},\mathbf{t}$；$\mu\gets3\max_ie_i/c^2$
\While{$\mu\ge1$}
  \For{$k=1$ to $K$}
    \State 对所有$i$：$e_i\gets\|R(\mathbf{q})\mathbf{p}_i+\mathbf{t}-\mathbf{y}_i\|_{W_i}^2$
    \State 对所有$i$：$w_i\gets\big(\mu c^2/(\mu c^2+e_i)\big)^2$
    \State $(\mathbf{q},\mathbf{t})\gets\textsc{QPEP-Global}\big(\sum_iw_ie_i\big)$
  \EndFor
  \State $\mu\gets\max(1,\mu/\gamma)$
\EndWhile
\State \Return $\mathbf{q},\mathbf{t},w,e$
\end{algorithmic}
\end{algorithm}

退火调度为$\mu_0\ge3\max_ie_i(z^{(0)})/c^2$，$\mu_{k+1}=\max(1,\mu_k/\gamma)$，
$\gamma\in[1.2,2.0]$，内层交替次数固定或自适应。$\mu_0$的选取使每个残差的\emph{径向}二阶导数
为正；它并不使完整代价关于$z$凸。

\subsection{不动点、曲率修正与精确黎曼Hessian}

\begin{lemma}[不动点$\Rightarrow$驻性]\label{lem:fixedpoint}
若$z^*$全局极小化$\sum_iw_i^*e_i$且$w_i^*=\rho'(e_i(z^*))$，则$(z^*,w^*)$是GNC不动点且
$\Phi(z^*)=0$。反之，$\Phi(z^*)=0$加上同样的加权全局极小性质给出不动点；仅有驻性不够。
\end{lemma}

\begin{lemma}[曲率修正$=$Schur补]\label{lem:schur}
在GNC不动点处，令$H_{\mathrm{IRLS}}=\sum_i\rho'(e_i^*)H_i$，
$\Delta H_{\mathrm{redesc}}=\sum_i\rho''(e_i^*)\,g_ig_i^\top\preceq0$，$D=[g_1,\dots,g_N]$，
$\Psi''=\diag(\psi''(w_i^*))$。联合BR目标的Hessian为
\begin{equation}
\begin{aligned}
H_J&=\begin{bmatrix}H_{\mathrm{IRLS}}&D\\D^\top&\Psi''\end{bmatrix},\\
H_{\mathrm{rob}}&=H_{\mathrm{IRLS}}-D\Psi''^{-1}D^\top
=H_{\mathrm{IRLS}}+\Delta H_{\mathrm{redesc}}.
\end{aligned}
\end{equation}
由于$\psi''(w_i^*)=-1/\rho''(e_i^*)>0$，重降修正正是被消去的权重块；且若
$H_{\mathrm{IRLS}}\succeq0$，则$H_J\succ0\iff H_{\mathrm{rob}}\succ0$。冻结权重会删除该
Schur补；第~\ref{sec:gradient}节量化由此产生的梯度误差。
\end{lemma}

\begin{lemma}[$\mathbb{S}^3\times\mathbb{R}^3$上的精确黎曼Hessian]\label{lem:hessian}
令$f\in C^2(\mathcal{M})$，欧氏梯度为$G$、Hessian为$H_E$。则
$\operatorname{Hess}f=\Pi\big(H_E-(\mathbf{q}^\top G_q)\diag(I_4,0_3)\big)\Pi$，且在正交切基$B$下
\begin{equation}
H_f=B^\top H_E B-(\mathbf{q}^\top G_q)\diag(I_3,0_3).
\end{equation}
对$F_{\mathrm{rob}}$，
\begin{equation}
\boxed{
\begin{aligned}
H_{\mathrm{rob}}&=B^\top\Big[\sum_i\rho'(e_i)H_i^E\Big]B\\
&\quad+B^\top\Big[\sum_i\rho''(e_i)g_i^E(g_i^E)^\top\Big]B\\
&\quad-\lambda_{\mathrm{rob}}\diag(I_3,0_3),\\
\lambda_{\mathrm{rob}}&=\mathbf{q}^\top G_q=\sum_i\rho'(e_i)\lambda_i .
\end{aligned}}
\end{equation}
若$H_i$与$g_i$表示黎曼量，则等价形式为
$H_{\mathrm{rob}}=\sum_i\rho'(e_i)H_i+\sum_i\rho''(e_i)g_ig_i^\top$。项
$-\lambda_{\mathrm{rob}}\diag(I_3,0_3)$是单位四元数约束的乘子，不能通过重缩放吸收：对3D--3D，
\begin{equation}
\begin{aligned}
\lambda_i&=\mathbf{q}^\top g_{i,q}^E=4\mathbf{u}_i^\top W_i\mathbf{r}_i\\
&=4\big(e_i-(\mathbf{t}-\mathbf{y}_i)^\top W_i\mathbf{r}_i\big)\neq4e_i
\end{aligned}
\end{equation}
只要$(\mathbf{t}-\mathbf{y}_i)^\top W_i\mathbf{r}_i\neq0$。在球面上，
$f(\mathbf{q})=\frac12\|\mathbf{q}\|^2$为常数、黎曼Hessian为零，而欧氏公式在切坐标下给出
$I_3$。在$c=1.0$的六维3D--3D例子中，未修正公式与测地有限差分Hessian的相对Frobenius范数差为
$25.9\%$，修正后公式吻合到$9.9\times10^{-8}$；乘子为$\lambda_{\mathrm{rob}}=1.033$。
\end{lemma}

\begin{lemma}[3D--3D权重]\label{lem:weights}
沿用上述记号，$g_{i,q}^E=2J_i^\top W_i\mathbf{r}_i$，$g_{i,t}^E=2W_i\mathbf{r}_i$，
$\lambda_i=4\mathbf{u}_i^\top W_i\mathbf{r}_i$。总乘子为
$\lambda_{\mathrm{rob}}=4\sum_i\rho'(e_i)\mathbf{u}_i^\top W_i\mathbf{r}_i$，而不是
$4\sum_i\rho'(e_i)e_i$。
\end{lemma}

\subsection{精确隐式微分}

\begin{theorem}[精确不动点IFT]\label{thm:ift}
在(A1)--(A4)下，设$\Phi(z^*,\theta_0)=0$且$H_{\mathrm{rob}}(z^*,\theta_0)$可逆。则存在$\theta_0$
的邻域$U$与唯一的$C^1$分支$\bar z:U\to\mathcal{M}$，满足$\bar z(\theta_0)=z^*$且
$\Phi(\bar z(\theta),\theta)=0$，并且
\begin{equation}
\frac{\mathrm{d}\bar z}{\mathrm{d}\theta}=-H_{\mathrm{rob}}^{-1}\frac{\partial\Phi}{\partial\theta}.
\end{equation}
若$H_{\mathrm{rob}}\succ0$，该分支是严格局部极小点，并在小扰动下保持；负特征值使$z^*$成为
鞍点；$H_{\mathrm{rob}}$奇异使隐函数定理失效，分支可能分岔。该导数是\emph{所选驻点分支}的导数；
若GNC迭代切换分支，导数对应其所选的分支。可逆性是IFT条件，而$\lambda_{\min}>0$是严格局部极小点
与分支持续性的更强条件。
\end{theorem}

\begin{theorem}[右端项分解]\label{thm:rhs}
对参数$\theta$，
\begin{equation}
\frac{\partial\Phi}{\partial\theta}
=\underbrace{\sum_i\rho'(e_i)\frac{\partial g_i}{\partial\theta}}_{\text{冻结权重项}}
+\underbrace{\sum_i\rho''(e_i)\,g_i\Big(\frac{\partial e_i}{\partial\theta}\Big)^\top}
_{\text{隐式权重敏感度}}.
\end{equation}
对3D--3D中的$\theta=\mathbf{y}_i$，
\begin{equation}
\begin{aligned}
\frac{\partial\Phi_t}{\partial\mathbf{y}_i}&=-2\big(\rho'(e_i)W_i
+2\rho''(e_i)W_i\mathbf{r}_i\mathbf{r}_i^\top W_i\big),\\
\frac{\partial\Phi_q}{\partial\mathbf{y}_i}&=-2J_i^\top W_i
\big(\rho'(e_i)I_3+2\rho''(e_i)\mathbf{r}_i\mathbf{r}_i^\top W_i\big).
\end{aligned}
\end{equation}
冻结权重近似只保留第一项，并令
$\partial\Phi_t^{\mathrm{naive}}/\partial\mathbf{y}_i=-2\rho'(e_i)W_i$。
\end{theorem}

\subsection{冻结权重梯度何时符号错误}
\label{sec:signflip}

取$W_i=I_3$，$\partial\Phi_t/\partial\mathbf{y}_i$的平行特征值为$-\lambda_\parallel(e_i)$，其中
\begin{equation}
\lambda_\parallel(e)=2(\rho'(e)+2e\rho''(e))=\frac{2c^4(c^2-3e)}{(c^2+e)^3}.
\end{equation}
于是$\lambda_\parallel(0)=2$，$\lambda_\parallel(c^2/3)=0$，$\lambda_\parallel(c^2)=-1/2$，且当
$e\to\infty$时$\lambda_\parallel(e)\sim-6c^4/e^2$。精确与朴素的平行特征值相差
$-4e\rho''=8c^4e/(c^2+e)^3>0$，故当$e>c^2/3$时二者符号严格相反；影响最大的残差是
$e=c^2$附近的\emph{边界}残差，而不是最远的粗大外点。

\begin{proposition}[一维符号翻转]\label{prop:signflip}
令$F(t)=\sum_i\rho((t-y_i)^2)$，核为GM。在满足
$H_{\mathrm{rob}}=\sum_i\lambda_\parallel(e_i)>0$的驻点$t^*$处，
\begin{equation}
\begin{aligned}
\frac{\mathrm{d}t^*}{\mathrm{d}y_j}&=\frac{\lambda_\parallel(e_j)}{H_{\mathrm{rob}}},\\
\Big(\frac{\mathrm{d}t^*}{\mathrm{d}y_j}\Big)_{\mathrm{naive}}
&=\frac{2\rho'(e_j)}{\sum_i2\rho'(e_i)}>0.
\end{aligned}
\end{equation}
若$e_j>c^2/3$，则$\lambda_\parallel(e_j)<0$，于是精确导数严格为负，而朴素导数严格为正。若
$H_{\mathrm{rob}}=0$，隐函数定理失效；若$H_{\mathrm{rob}}<0$，该点不是局部极小点。
\end{proposition}

\begin{proposition}[存在性]\label{prop:existence}
对任意$c>0$及充分大的$y_2-y_1$，两观测代价有严格局部极小点$t^*$，满足
$e_1<c^2/3<e_2$且$H_{\mathrm{rob}}>0$；因此命题~\ref{prop:signflip}的符号翻转确实发生。证明利用
$\Phi$在$0$处与$Kc^4/d^3$处的符号，以及极限$H_{\mathrm{rob}}\to\lambda_\parallel(0)=2>0$。
\end{proposition}

在6D中，位姿灵敏度的符号\emph{不}仅由$\lambda_\parallel$决定：
$\partial z^*/\partial\mathbf{y}_i=-H_{\mathrm{rob}}^{-1}\partial\Phi/\partial\mathbf{y}_i$，
而逆Hessian会混合各个方向。命题~\ref{prop:signflip}是严格陈述；6D对应为数值结果，余弦为
$-0.079$，6D类比需要对$H_{\mathrm{rob}}$与$\partial\Phi/\partial\mathbf{y}_i$的联合谱条件。

\begin{proposition}[逐观测$O(\|\mathbf{r}\|^{-4})$界]\label{prop:bound}
在(A1)、(A5)下，设沿路径$\|\mathbf{y}_i\|\to\infty$时$z^*(\mathbf{y}_i)$保持在紧集中，且
$H_{\mathrm{rob}}$在该处一致正定、$\lambda_{\min}\ge\mu_0>0$。则
\begin{equation}
\begin{aligned}
\Big\|\frac{\partial\Phi}{\partial\mathbf{y}_i}\Big\|_F&\le C\frac{c^4}{(c^2+e_i)^2}
\le C\frac{c^4}{\|\mathbf{r}_i\|^4},\\
\Big\|\frac{\partial z^*}{\partial\mathbf{y}_i}\Big\|_F&\le\frac{C}{\mu_0}
\frac{c^4}{\|\mathbf{r}_i\|^4}.
\end{aligned}
\end{equation}
前向影响仅为$O(\|\mathbf{r}_i\|^{-3})$：
$\|\rho'(e_i)g_i\|\le2\overline wc^4/\|\mathbf{r}_i\|^3$。该界是逐观测的；它不能阻止$N$个外点
累积，也不约束成为全局极小的相干外点簇。因此单个观测上的梯度衰减与前向鲁棒性是两个不同性质。
\end{proposition}

\subsection{算法、复杂度与二阶认证}

反向传播（算法~\ref{alg:backward}，附录~\ref{app:backward}）存储$z^*,w^*$并重算
$g_i^E,H_i^E$；它对每个参数块求解一个$6\times6$系统，且从不展开GNC迭代。内存为权重与残差的
$O(N)$、GNC迭代次数的$O(1)$；复杂度为每次前向迭代$O(N)$加一次固定规模的代数求解，反向装配为
$O(N)$加$O(6^3)$。\textbf{二阶认证}是$H_{\mathrm{rob}}$的谱：$H_{\mathrm{rob}}\succ0$使不动点
成为联合BR目标的严格局部极小点，负特征值使其成为鞍点，零穿越使分支分岔且梯度范数爆炸。这是局部
稳定性诊断——一个错误的、与外点一致的位姿也可以是$H_{\mathrm{rob}}\succ0$的严格局部极小点。
